\section{数据驱动方法}

\subsection{从规则驱动到数据驱动}

传统计算机科学中的许多问题（如排序）可以用明确的算法流程（if-then-else + 循环）解决。但图像分类无法这样做——我们无法为每种物体的每种外观编写规则。

\begin{figure}[H]
\centering
\includegraphics[width=0.85\textwidth]{slides-images/slide-017.jpg}
\caption{早期尝试：基于边缘检测和角点特征的规则方法。Hubel \& Wiesel 的研究启发了通过边缘、角点等低级特征识别物体的思路，但难以扩展\protect\footnotemark}
\end{figure}
\footnotetext{Slides 第18页。}

早期曾有研究者尝试通过边缘检测器提取边缘、识别角点模式，然后基于这些特征建立分类规则。虽然在有限场景下取得了一些成功，但存在两个根本问题：（1）每种物体都需要单独设计识别规则，\textbf{无法扩展}；（2）规则的设计本身就极其困难。

\subsection{数据驱动方法的三步流程}

机器学习提供了一种全新的范式——\textbf{数据驱动方法}（Data-Driven Approach）：

\begin{figure}[H]
\centering
\includegraphics[width=0.85\textwidth]{slides-images/slide-019.jpg}
\caption{数据驱动方法的三步流程\protect\footnotemark}
\end{figure}
\footnotetext{Slides 第20页。}

\begin{importantbox}{数据驱动方法三步流程}
\begin{enumerate}
    \item \textbf{收集数据}：收集大量图像及其对应的标签
    \item \textbf{训练分类器}：使用机器学习算法从数据中学习一个模型——实现 \texttt{train(images, labels) -> model}
    \item \textbf{评估分类器}：在新的、未见过的图像上评估模型性能——实现 \texttt{predict(model, test\_images) -> labels}
\end{enumerate}
\end{importantbox}

\begin{figure}[H]
\centering
\includegraphics[width=0.85\textwidth]{slides-images/slide-020.jpg}
\caption{数据驱动方法的代码框架：\texttt{train} 函数学习模型，\texttt{predict} 函数使用模型预测新图像的标签\protect\footnotemark}
\end{figure}
\footnotetext{Slides 第21页。}

这种方法的核心优势在于：不需要为每种物体手工设计规则，而是让算法从数据中自动学习区分不同类别的模式。

\subsection{本章小结}

数据驱动方法通过``收集数据 $\rightarrow$ 训练模型 $\rightarrow$ 评估预测''的三步流程，替代了传统的手工规则方法。这一范式是现代计算机视觉和深度学习的基础。接下来我们将介绍两种具体的数据驱动分类器。

%% ========== Part 3: K 近邻分类器 ==========

